% Physical Review Letters (PRL) Paper Template
% Generated by ARIS paper-write skill
% Uses RevTeX4-2 (included in standard TeX distributions: texlive, miktex)

\documentclass[reprint,aps,prl,superscriptaddress,showpacs,floatfix]{revtex4-2}
% Options:
%   reprint  — single-column reprint format (use "twocolumn" for final check)
%   aps      — American Physical Society
%   prl      — Physical Review Letters style
%   superscriptaddress — compact author/affiliation format
%   showpacs — show PACS numbers (optional)
%   floatfix — better float placement

% Math
\usepackage{amsmath,amssymb,amsfonts,mathtools}

% Typography & Graphics
\usepackage{graphicx}
\usepackage{xcolor}
\usepackage{booktabs}
\usepackage{subcaption}
\usepackage{bm}            % bold math symbols (common in physics)
\IfFileExists{siunitx.sty}{\usepackage{siunitx}}{} % SI units (common in physics, optional)

% hyperref — RevTeX loads natbib internally, so load hyperref carefully
\usepackage{hyperref}
\hypersetup{
    colorlinks=true,
    linkcolor=blue,
    citecolor=blue,
    urlcolor=blue
}

% cleveref AFTER hyperref
\usepackage[capitalize,noabbrev]{cleveref}

% Shared math commands
\input{math_commands}

\begin{document}

\preprint{APS/123-QED}  % optional preprint number

\title{Paper Title Here}

%% Authors — PRL is NOT anonymous
\author{First Author}
\affiliation{Department of Physics, University One, City, Country}

\author{Second Author}
\affiliation{Department of Physics, University One, City, Country}
\affiliation{Department of Physics, University Two, City, Country}

\author{Corresponding Author}
\email{corresponding@example.com}
\affiliation{Department of Physics, University Two, City, Country}

\date{\today}

\begin{abstract}
\input{sections/0_abstract}
\end{abstract}

% \pacs{xx.xx.Xx}  % Optional: PACS numbers

\maketitle

%% PRL is very compact (4 journal pages ≈ 3750 words).
%% Sections are optional — many PRL papers use them sparingly.
%% Common pattern: Introduction (no header) → Method → Results → Conclusion
%% Or just continuous text with minimal section breaks.

\input{sections/1_introduction}
\input{sections/2_method}
\input{sections/3_results}
\input{sections/4_conclusion}

\begin{acknowledgments}
We thank X for helpful discussions.
% Funding: This work was supported by ...
\end{acknowledgments}

\bibliography{references}

\end{document}

%% === Supplemental Material ===
%% PRL supplemental material is a SEPARATE file (supplement.tex).
%% It is NOT appended to the main paper.
%% Create a separate supplement.tex if needed:
%%   \documentclass[aps,prl,reprint,superscriptaddress]{revtex4-2}
%%   \begin{document}
%%   \title{Supplemental Material for "Paper Title"}
%%   \maketitle
%%   ... content ...
%%   \end{document}
